%  LaTeX support: latex@mdpi.com 
%  For support, please attach all files needed for compiling as well as the log file, and specify your operating system, LaTeX version, and LaTeX editor.

%=================================================================
%\documentclass[journal,article,submit,pdftex,moreauthors]{Definitions/mdpi} 
\documentclass[preprints,article,submit,pdftex,moreauthors]{Definitions/mdpi} 
% For posting an early version of this manuscript as a preprint, you may use "preprints" as the journal. Changing "submit" to "accept" before posting will remove line numbers.

% Below journals will use APA reference format:
% admsci, behavsci, businesses, econometrics, economies, education, ejihpe, games, humans, ijfs, journalmedia, jrfm, languages, psycholint, publications, tourismhosp, youth

% Below journals will use Chicago reference format:
% arts, genealogy, histories, humanities, jintelligence, laws, literature, religions, risks, socsci

%--------------------
% Class Options:
%--------------------
%----------
% journal
%----------
% Choose between the following MDPI journals:
% accountaudit, acoustics, actuators, addictions, adhesives, admsci, adolescents, aerobiology, aerospace, agriculture, agriengineering, agrochemicals, agronomy, ai, air, algorithms, allergies, alloys, amh, analytica, analytics, anatomia, anesthres, animals, antibiotics, antibodies, antioxidants, applbiosci, appliedchem, appliedmath, appliedphys, applmech, applmicrobiol, applnano, applsci, aquacj, architecture, arm, arthropoda, arts, asc, asi, astronomy, atmosphere, atoms, audiolres, automation, axioms, bacteria, batteries, bdcc, behavsci, beverages, biochem, bioengineering, biologics, biology, biomass, biomechanics, biomed, biomedicines, biomedinformatics, biomimetics, biomolecules, biophysica, biosensors, biosphere, biotech, birds, blockchains, bloods, blsf, brainsci, breath, buildings, businesses, cancers, carbon, cardiogenetics, catalysts, cells, ceramics, challenges, chemengineering, chemistry, chemosensors, chemproc, children, chips, cimb, civileng, cleantechnol, climate, clinbioenerg, clinpract, clockssleep, cmd, cmtr, coasts, coatings, colloids, colorants, commodities, complications, compounds, computation, computers, condensedmatter, conservation, constrmater, cosmetics, covid, crops, cryo, cryptography, crystals, csmf, ctn, curroncol, cyber, dairy, data, ddc, dentistry, dermato, dermatopathology, designs, devices, diabetology, diagnostics, dietetics, digital, disabilities, diseases, diversity, dna, drones, dynamics, earth, ebj, ecm, ecologies, econometrics, economies, education, eesp, ejihpe, electricity, electrochem, electronicmat, electronics, encyclopedia, endocrines, energies, eng, engproc, ent, entomology, entropy, environments, epidemiologia, epigenomes, esa, est, famsci, fermentation, fibers, fintech, fire, fishes, fluids, foods, forecasting, forensicsci, forests, fossstud, foundations, fractalfract, fuels, future, futureinternet, futureparasites, futurepharmacol, futurephys, futuretransp, galaxies, games, gases, gastroent, gastrointestdisord, gastronomy, gels, genealogy, genes, geographies, geohazards, geomatics, geometry, geosciences, geotechnics, geriatrics, glacies, grasses, greenhealth, gucdd, hardware, hazardousmatters, healthcare, hearts, hemato, hematolrep, heritage, higheredu, highthroughput, histories, horticulturae, hospitals, humanities, humans, hydrobiology, hydrogen, hydrology, hygiene, idr, iic, ijerph, ijfs, ijgi, ijmd, ijms, ijns, ijpb, ijt, ijtm, ijtpp, ime, immuno, informatics, information, infrastructures, inorganics, insects, instruments, inventions, iot, j, jal, jcdd, jcm, jcp, jcs, jcto, jdad, jdb, jeta, jfb, jfmk, jimaging, jintelligence, jlpea, jmahp, jmmp, jmms, jmp, jmse, jne, jnt, jof, joitmc, joma, jop, jor, journalmedia, jox, jpbi, jpm, jrfm, jsan, jtaer, jvd, jzbg, kidney, kidneydial, kinasesphosphatases, knowledge, labmed, laboratories, land, languages, laws, life, lights, limnolrev, lipidology, liquids, literature, livers, logics, logistics, lubricants, lymphatics, machines, macromol, magnetism, magnetochemistry, make, marinedrugs, materials, materproc, mathematics, mca, measurements, medicina, medicines, medsci, membranes, merits, metabolites, metals, meteorology, methane, metrics, metrology, micro, microarrays, microbiolres, microelectronics, micromachines, microorganisms, microplastics, microwave, minerals, mining, mmphys, modelling, molbank, molecules, mps, msf, mti, multimedia, muscles, nanoenergyadv, nanomanufacturing, nanomaterials, ncrna, ndt, network, neuroglia, neurolint, neurosci, nitrogen, notspecified, nursrep, nutraceuticals, nutrients, obesities, oceans, ohbm, onco, oncopathology, optics, oral, organics, organoids, osteology, oxygen, parasites, parasitologia, particles, pathogens, pathophysiology, pediatrrep, pets, pharmaceuticals, pharmaceutics, pharmacoepidemiology, pharmacy, philosophies, photochem, photonics, phycology, physchem, physics, physiologia, plants, plasma, platforms, pollutants, polymers, polysaccharides, populations, poultry, powders, preprints, proceedings, processes, prosthesis, proteomes, psf, psych, psychiatryint, psychoactives, psycholint, publications, purification, quantumrep, quaternary, qubs, radiation, reactions, realestate, receptors, recycling, regeneration, religions, remotesensing, reports, reprodmed, resources, rheumato, risks, robotics, rsee, ruminants, safety, sci, scipharm, sclerosis, seeds, sensors, separations, sexes, signals, sinusitis, siuj, skins, smartcities, sna, societies, socsci, software, soilsystems, solar, solids, spectroscj, sports, standards, stats, std, stresses, surfaces, surgeries, suschem, sustainability, symmetry, synbio, systems, tae, targets, taxonomy, technologies, telecom, test, textiles, thalassrep, therapeutics, thermo, timespace, tomography, tourismhosp, toxics, toxins, transplantology, transportation, traumacare, traumas, tropicalmed, universe, urbansci, uro, vaccines, vehicles, venereology, vetsci, vibration, virtualworlds, viruses, vision, waste, water, wem, wevj, wild, wind, women, world, youth, zoonoticdis

%---------
% article
%---------
% The default type of manuscript is "article", but can be replaced by: 
% abstract, addendum, article, benchmark, book, bookreview, briefcommunication, briefreport, casereport, changes, clinicopathologicalchallenge, comment, commentary, communication, conceptpaper, conferenceproceedings, correction, conferencereport, creative, datadescriptor, discussion, entry, expressionofconcern, extendedabstract, editorial, essay, erratum, fieldguide, hypothesis, interestingimages, letter, meetingreport, monograph, newbookreceived, obituary, opinion, proceedingpaper, projectreport, reply, retraction, review, perspective, protocol, shortnote, studyprotocol, supfile, systematicreview, technicalnote, viewpoint, guidelines, registeredreport, tutorial,  giantsinurology, urologyaroundtheworld
% supfile = supplementary materials

%----------
% submit
%----------
% The class option "submit" will be changed to "accept" by the Editorial Office when the paper is accepted. This will only make changes to the frontpage (e.g., the logo of the journal will get visible), the headings, and the copyright information. Also, line numbering will be removed. Journal info and pagination for accepted papers will also be assigned by the Editorial Office.

%------------------
% moreauthors
%------------------
% If there is only one author the class option oneauthor should be used. Otherwise use the class option moreauthors.

%---------
% pdftex
%---------
% The option pdftex is for use with pdfLaTeX. Remove "pdftex" for (1) compiling with LaTeX & dvi2pdf (if eps figures are used) or for (2) compiling with XeLaTeX.

%=================================================================
% MDPI internal commands - do not modify
\firstpage{1} 
\makeatletter 
\setcounter{page}{\@firstpage} 
\makeatother
\pubvolume{1}
\issuenum{1}
\articlenumber{0}
\pubyear{2026}
\copyrightyear{2026}
%\externaleditor{Firstname Lastname} % More than 1 editor, please add `` and '' before the last editor name
\datereceived{ } 
\daterevised{ } % Comment out if no revised date
\dateaccepted{ } 
\datepublished{ } 
%\datecorrected{} % For corrected papers: "Corrected: XXX" date in the original paper.
%\dateretracted{} % For retracted papers: "Retracted: XXX" date in the original paper.
\hreflink{https://doi.org/} % If needed use \linebreak
%\doinum{}
%\pdfoutput=1 % Uncommented for upload to arXiv.org
%\CorrStatement{yes}  % For updates
%\longauthorlist{yes} % For many authors that exceed the left citation part

%=================================================================
% Add packages and commands here. The following packages are loaded in our class file: fontenc, inputenc, calc, indentfirst, fancyhdr, graphicx, epstopdf, lastpage, ifthen, float, amsmath, amssymb, lineno, setspace, enumitem, mathpazo, booktabs, titlesec, etoolbox, tabto, xcolor, colortbl, soul, multirow, microtype, tikz, totcount, changepage, attrib, upgreek, array, tabularx, pbox, ragged2e, tocloft, marginnote, marginfix, enotez, amsthm, natbib, hyperref, cleveref, scrextend, url, geometry, newfloat, caption, draftwatermark, seqsplit
% cleveref: load \crefname definitions after \begin{document}

%=================================================================
% Please use the following mathematics environments: Theorem, Lemma, Corollary, Proposition, Characterization, Property, Problem, Example, ExamplesandDefinitions, Hypothesis, Remark, Definition, Notation, Assumption
%% For proofs, please use the proof environment (the amsthm package is loaded by the MDPI class).

%=================================================================
% Full title of the paper (Capitalized)
\Title{PARCS-Agent: Enabling AI Agents to Exploit Distributed Parallel Computing via the Model Context Protocol}

% MDPI internal command: Title for citation in the left column
\TitleCitation{PARCS-Agent: AI Agents and Distributed Parallel Computing via MCP}

% Author Orchid ID: enter ID or remove command
\newcommand{\orcidauthorA}{0000-0000-0000-000X}

% Authors, for the paper (add full first names)
\Author{Oleksii Bohusevych $^{1}$\orcidA{}}

% MDPI internal command: Authors, for metadata in PDF
\AuthorNames{Oleksii Bohusevych}

\isAPAStyle{%
       \AuthorCitation{Bohusevych, O.}
         }{%
        \isChicagoStyle{%
        \AuthorCitation{Bohusevych, Oleksii.}
        }{
        \AuthorCitation{Bohusevych, O.}
        }
}

% Affiliations / Addresses
\address{%
$^{1}$ \quad Faculty of Computer Science and Cybernetics, Taras Shevchenko National University of Kyiv, 01601 Kyiv, Ukraine; oleksii.bohusevych@knu.ua}

% Contact information of the corresponding author
\corres{Correspondence: oleksii.bohusevych@knu.ua}

% Abstract
\abstract{%
Large language model (LLM) agents have demonstrated impressive capabilities in tool use, code generation, and scientific reasoning, yet they remain bound by the throughput of a single execution thread. Computationally intensive tasks—Monte Carlo simulation, combinatorial optimisation, hyperparameter search, large-scale graph analysis—either time out or require severe undersampling when handled sequentially. We present PARCS-Agent, a system that exposes a cloud-scale distributed computing cluster to an AI agent through three Model Context Protocol (MCP) tools: \texttt{get\_cluster\_info}, \texttt{create\_session}, and \texttt{run\_layer}. The agent submits arbitrary C\# computation logic; the platform compiles it with Roslyn at runtime, fans it out across up to 21 Kubernetes worker pods via Google Cloud Pub/Sub and KEDA autoscaling, and returns aggregated results. We introduce the PARCS-Agent Benchmark, a suite of 15 tasks spanning quantitative finance, bioinformatics, epidemiology, machine learning, cryptography, and atmospheric science, designed so that the sequential wall-clock cost substantially exceeds practical agent time budgets. Experiments show that PARCS-enabled agents achieve mean speedups of 15--18$\times$ over sequential baselines while delivering statistically equivalent or superior solution quality. PARCS-Agent is, to our knowledge, the first system to use MCP as the interface between an AI agent and a live distributed HPC cluster, and the first benchmark to measure \emph{speedup ratio} as a first-class agent capability metric.%
}

% Keywords
\keyword{AI agents; Model Context Protocol; distributed computing; parallel computation; LLM tool use; benchmarking; Kubernetes; high-performance computing} 

% The fields PACS, MSC, and JEL may be left empty or commented out if not applicable
%\PACS{J0101}
%\MSC{}
%\JEL{}

%%%%%%%%%%%%%%%%%%%%%%%%%%%%%%%%%%%%%%%%%%
% Only for the journal Diversity
%\LSID{\url{http://}}

%%%%%%%%%%%%%%%%%%%%%%%%%%%%%%%%%%%%%%%%%%
% Only for the journal Applied Sciences
%\featuredapplication{Authors are encouraged to provide a concise description of the specific application or a potential application of the work. This section is not mandatory.}
%%%%%%%%%%%%%%%%%%%%%%%%%%%%%%%%%%%%%%%%%%

%%%%%%%%%%%%%%%%%%%%%%%%%%%%%%%%%%%%%%%%%%
% Only for the journal Data
%\dataset{DOI number or link to the deposited data set if the data set is published separately. If the data set shall be published as a supplement to this paper, this field will be filled by the journal editors. In this case, please submit the data set as a supplement.}
%\datasetlicense{License under which the data set is made available (CC0, CC-BY, CC-BY-SA, CC-BY-NC, etc.)}

%%%%%%%%%%%%%%%%%%%%%%%%%%%%%%%%%%%%%%%%%%
% Only for the journal Toxins
%\keycontribution{The breakthroughs or highlights of the manuscript. Authors can write one or two sentences to describe the most important part of the paper.}

%%%%%%%%%%%%%%%%%%%%%%%%%%%%%%%%%%%%%%%%%%
% Only for the journal Encyclopedia
%\encyclopediadef{For entry manuscripts only: please provide a brief overview of the entry title instead of an abstract.}

%%%%%%%%%%%%%%%%%%%%%%%%%%%%%%%%%%%%%%%%%%
% Only for the journal Advances in Respiratory Medicine and Smart Cities
%\addhighlights{yes}
%\renewcommand{\addhighlights}{%

%\noindent This is an obligatory section in “Advances in Respiratory Medicine'' and ``Smart Cities”, whose goal is to increase the discoverability and readability of the article via search engines and other scholars. Highlights should not be a copy of the abstract, but a simple text allowing the reader to quickly and simplified find out what the article is about and what can be cited from it. Each of these parts should be devoted up to 2~bullet points.\vspace{3pt}\\
%\textbf{What are the main findings?}
% \begin{itemize}[labelsep=2.5mm,topsep=-3pt]
% \item First bullet.
% \item Second bullet.
% \end{itemize}\vspace{3pt}
%\textbf{What are the implications of the main findings?}
% \begin{itemize}[labelsep=2.5mm,topsep=-3pt]
% \item First bullet.
% \item Second bullet.
% \end{itemize}
%}

%%%%%%%%%%%%%%%%%%%%%%%%%%%%%%%%%%%%%%%%%%
\begin{document}

%%%%%%%%%%%%%%%%%%%%%%%%%%%%%%%%%%%%%%%%%%
%%%%%%%%%%%%%%%%%%%%%%%%%%%%%%%%%%%%%%%%%%
\section{Introduction}
% TODO: to be written

%%%%%%%%%%%%%%%%%%%%%%%%%%%%%%%%%%%%%%%%%%
\section{Related Work}
\label{sec:related}

\subsection{Tool-Augmented Language Model Agents}

Two main directions can be distinguished in the literature on augmenting language models with external tools. The first, introduced by Toolformer by Schick et al.~\citep{ref-toolformer}, showed that a model is capable of learning, without supervised examples, when and how to call external APIs such as calculators, search engines, or translation services. The second, ReAct by Yao et al.~\citep{ref-react}, proposed a tighter integration of reasoning and action: the model produces a chain-of-thought trace, takes an action, receives an observation, and repeats. Both works established the agent-tool loop as a general paradigm, and many subsequent systems build directly on top of them.

The scope of tool use has since expanded considerably. ToolLLM by Qin et al.~\citep{ref-toolllm} curated a collection of 16{,}000 real-world REST endpoints and trained a model to select and invoke them using a tree-based search strategy. Gorilla by Patil et al.~\citep{ref-gorilla} addressed the specific problem of hallucinated API signatures when working with ML framework APIs, showing that retrieval-augmented generation can mitigate this. SciAgent by Ma et al.~\citep{ref-sciagent} brought scientific computing tools into the picture, assembling over 30{,}000 tool-annotated examples spanning physics, chemistry, mathematics, and finance, alongside a dedicated benchmark. Among the broader surveys, the one by Mialon et al.~\citep{ref-augmented-survey} is notable for placing all of the above within a taxonomy of memory, reasoning, and action, and for noting that large-scale computation as a tool type remains underexplored.

In all of the aforementioned systems, a tool call is stateless and local: it takes inputs, returns a result, and exits. PARCS-Agent departs from this model, as its \texttt{run\_layer} tool submits user-written C\# code to a running Kubernetes cluster, fans it out across a pool of workers via a Pub/Sub queue, and aggregates their output. Hence, a single tool invocation may correspond to dozens of parallel processes rather than a single function call.

LLMCompiler by Kim et al.~\citep{ref-llmcompiler} is worth singling out, as it does introduce parallelism into tool use. It analyses the dependency structure of a query, builds a DAG of tool calls, and dispatches independent branches concurrently, achieving up to $3.7\times$ latency reduction on multi-hop tasks. The said parallelism, however, applies to independent calls to existing tools, not to distributed execution of user-supplied code. CodeAct by Wang et al.~\citep{ref-codeact} proposes Python as the agent's action space and runs it in a stateful interpreter that persists across turns. At the same time, execution remains single-threaded and local, making it unsuitable for computationally intensive tasks.

\subsection{LLMs for High-Performance and Parallel Computing}

On the one hand, there is a growing body of work examining whether LLMs can produce correct parallel code. ParEval by Nichols et al.~\citep{ref-pareval} comprises 420 problems drawn from numerical computing, graph algorithms, and Monte Carlo methods, and evaluates generated code targeting MPI, OpenMP, CUDA, and Kokkos. The findings indicate that state-of-the-art models generate significantly less correct parallel code than serial code for the same problems, with correctness declining notably as task complexity grows. HPC-Coder-V2~\citep{ref-hpccoder} brought improvements in this direction by fine-tuning on a corpus of 120{,}000 parallel code samples, reporting gains on MPI and OpenMP generation in particular. Separately, recent frontier models were evaluated on classical HPC kernels such as matrix multiplication and conjugate gradient, with results indicating that non-trivial memory access patterns remain problematic~\citep{ref-deepseekhpc}.

On the other hand, PARCS-Agent sidesteps the said code generation problem altogether. The agent writes only per-worker logic as a plain C\# class implementing a fixed interface; the platform compiles it via Roslyn, distributes the resulting binary to worker pods, and handles all coordination. As such, the agent is never required to reason about parallelism in code, and the gap identified by ParEval becomes irrelevant when the parallel infrastructure is a managed service.

He et al.~\citep{ref-legionllm} applied LLMs to a different HPC problem, namely the automated optimisation of the task-to-processor mapper in the Legion distributed runtime. The Legion mapper DSL served as the target representation, and the system proposed mapping strategies that previously required days of expert tuning. It is worth mentioning, however, that the said work presupposes an already-parallelised application; the problem it addresses is performance tuning rather than enabling new computations to be executed in parallel from scratch.

\subsection{Agent Benchmarks for Computational Tasks}

Several benchmarks have been proposed in recent years to measure agent performance on complex, multi-step tasks. GAIA~\citep{ref-gaia} comprises 466 questions requiring a combination of web browsing, file parsing, and multi-modal reasoning; the hardest third of questions remains largely unsolved. SWE-bench~\citep{ref-swebench} evaluates whether agents can resolve real GitHub issues by producing patches that pass the repository's test suite, with top systems solving roughly half of the verified subset. ScienceAgentBench~\citep{ref-sciagentbench} consists of 102 tasks drawn from 44 published papers, spanning bioinformatics, chemistry, and geographical information science, with the best agent completing around one-third autonomously. TheAgentCompany~\citep{ref-agentcompany} takes a different angle, benchmarking agents on workplace tasks in a simulated company environment.

All of the above benchmarks share a common assumption: tasks are designed to be solvable by a single-threaded agent within a reasonable wall-clock budget, and compute time is not a measured variable. The PARCS-Agent Benchmark is constructed around the opposite premise. Each of the 15 tasks carries a sequential cost that is impractical to complete within a reasonable budget, and the central result is the ratio of wall-clock time between a PARCS-enabled agent and a sequential baseline. Hence, speedup ratio serves as a primary evaluation criterion, which, to the best of the author's knowledge, has not been used in agent benchmarking before.

\subsection{The Model Context Protocol}

The Model Context Protocol (MCP) was introduced by Anthropic in 2024 as a standard interface for connecting AI agents to external tools, data, and prompts via JSON-RPC over server-sent events~\citep{ref-mcp}. Hou et al.~\citep{ref-mcpsurvey} surveyed the MCP ecosystem and identified security concerns comprising prompt injection, tool poisoning, and privilege escalation; notably, no prior work on exposing distributed compute infrastructure via MCP was found. A separate analysis of published MCP tool descriptions found that over half lacked adequate documentation of purpose or parameter constraints~\citep{ref-mcpsmell}. The said findings motivated a deliberate effort in the design of the PARCS-Agent tools, each of which specifies parameter bounds, error conditions, and expected return values explicitly.

As far as the author is aware, PARCS-Agent is the first system to expose a live parallel computing cluster as a set of MCP tools, and the first benchmark to evaluate agent performance specifically on tasks where access to such a cluster is the deciding factor.

%%%%%%%%%%%%%%%%%%%%%%%%%%%%%%%%%%%%%%%%%%
\section{System Architecture}
% TODO: to be written

%%%%%%%%%%%%%%%%%%%%%%%%%%%%%%%%%%%%%%%%%%
\section{PARCS-Agent Benchmark}
% TODO: to be written

%%%%%%%%%%%%%%%%%%%%%%%%%%%%%%%%%%%%%%%%%%
\section{Experiments}
% TODO: to be written

%%%%%%%%%%%%%%%%%%%%%%%%%%%%%%%%%%%%%%%%%%
\section{Discussion}
% TODO: to be written

%%%%%%%%%%%%%%%%%%%%%%%%%%%%%%%%%%%%%%%%%%
\section{Conclusions}
% TODO: to be written

%%%%%%%%%%%%%%%%%%%%%%%%%%%%%%%%%%%%%%%%%%
\vspace{6pt} 

%%%%%%%%%%%%%%%%%%%%%%%%%%%%%%%%%%%%%%%%%%
%% optional
%\supplementary{The following supporting information can be downloaded at:  \linksupplementary{s1}, Figure S1: title; Table S1: title; Video S1: title.}

% Only for journal Methods and Protocols:
% If you wish to submit a video article, please do so with any other supplementary material.
% \supplementary{The following supporting information can be downloaded at: \linksupplementary{s1}, Figure S1: title; Table S1: title; Video S1: title. A supporting video article is available at doi: link.}

% Only used for preprtints:
% \supplementary{The following supporting information can be downloaded at the website of this paper posted on \href{https://www.preprints.org/}{Preprints.org}.}

% Only for journal Hardware:
% If you wish to submit a video article, please do so with any other supplementary material.
% \supplementary{The following supporting information can be downloaded at: \linksupplementary{s1}, Figure S1: title; Table S1: title; Video S1: title.\vspace{6pt}\\
%\begin{tabularx}{\textwidth}{lll}
%\toprule
%\textbf{Name} & \textbf{Type} & \textbf{Description} \\
%\midrule
%S1 & Python script (.py) & Script of python source code used in XX \\
%S2 & Text (.txt) & Script of modelling code used to make Figure X \\
%S3 & Text (.txt) & Raw data from experiment X \\
%S4 & Video (.mp4) & Video demonstrating the hardware in use \\
%... & ... & ... \\
%\bottomrule
%\end{tabularx}
%}

%%%%%%%%%%%%%%%%%%%%%%%%%%%%%%%%%%%%%%%%%%
\authorcontributions{Conceptualization, O.B.; methodology, O.B.; software, O.B.; validation, O.B.; formal analysis, O.B.; investigation, O.B.; data curation, O.B.; writing - original draft preparation, O.B. The author has read and agreed to the published version of the manuscript.}

\funding{This research received no external funding.}

\institutionalreview{Not applicable.}

\informedconsent{Not applicable.}

\dataavailability{The benchmark dataset is publicly available at \url{https://huggingface.co/datasets/parcs-benchmark/parcs-agent-benchmark}. The platform source code is available at \url{https://github.com/alexbohuslav/parcs7}.} 

% Only for journal Drones
%\durcstatement{Current research is limited to the [please insert a specific academic field, e.g., XXX], which is beneficial [share benefits and/or primary use] and does not pose a threat to public health or national security. Authors acknowledge the dual-use potential of the research involving xxx and confirm that all necessary precautions have been taken to prevent potential misuse. As an ethical responsibility, authors strictly adhere to relevant national and international laws about DURC. Authors advocate for responsible deployment, ethical considerations, regulatory compliance, and transparent reporting to mitigate misuse risks and foster beneficial outcomes.}

% Only for journal Nursing Reports
%\publicinvolvement{Please describe how the public (patients, consumers, carers) were involved in the research. Consider reporting against the GRIPP2 (Guidance for Reporting Involvement of Patients and the Public) checklist. If the public were not involved in any aspect of the research add: ``No public involvement in any aspect of this research''.}
%
%% Only for journal Nursing Reports
%\guidelinesstandards{Please add a statement indicating which reporting guideline was used when drafting the report. For example, ``This manuscript was drafted against the XXX (the full name of reporting guidelines and citation) for XXX (type of research) research''. A complete list of reporting guidelines can be accessed via the equator network: \url{https://www.equator-network.org/}.}
%
%% Only for journal Nursing Reports
%\useofartificialintelligence{Please describe in detail any and all uses of artificial intelligence (AI) or AI-assisted tools used in the preparation of the manuscript. This may include, but is not limited to, language translation, language editing and grammar, or generating text. Alternatively, please state that “AI or AI-assisted tools were not used in drafting any aspect of this manuscript”.}

\acknowledgments{In this section you can acknowledge any support given which is not covered by the author contribution or funding sections. This may include administrative and technical support, or donations in kind (e.g., materials used for experiments).}

\conflictsofinterest{Declare conflicts of interest or state ``The authors declare no conflicts of interest.'' Authors must identify and declare any personal circumstances or interest that may be perceived as inappropriately influencing the representation or interpretation of reported research results. Any role of the funders in the design of the study; in the collection, analyses or interpretation of data; in the writing of the manuscript; or in the decision to publish the results must be declared in this section. If there is no role, please state ``The funders had no role in the design of the study; in the collection, analyses, or interpretation of data; in the writing of the manuscript; or in the decision to publish the results''.} 

%%%%%%%%%%%%%%%%%%%%%%%%%%%%%%%%%%%%%%%%%%
%% Optional

%% Only for journal Encyclopedia
%\entrylink{The Link to this entry published on the encyclopedia platform.}

\abbreviations{Abbreviations}{
The following abbreviations are used in this manuscript:\\

\noindent
\begin{tabular}{@{}ll}
MCP  & Model Context Protocol\\
LLM  & Large Language Model\\
HPC  & High-Performance Computing\\
API  & Application Programming Interface\\
KEDA & Kubernetes Event-Driven Autoscaling\\
GKE  & Google Kubernetes Engine\\
SA   & Simulated Annealing\\
VaR  & Value-at-Risk\\
CVaR & Conditional Value-at-Risk
\end{tabular}
}

%%%%%%%%%%%%%%%%%%%%%%%%%%%%%%%%%%%%%%%%%%
%% Optional
\appendixtitles{no} % Leave argument "no" if all appendix headings stay EMPTY (then no dot is printed after "Appendix A"). If the appendix sections contain a heading then change the argument to "yes".
\appendixstart
\appendix
\section[\appendixname~\thesection]{}
\subsection[\appendixname~\thesubsection]{}
The appendix is an optional section that can contain details and data supplemental to the main text---for example, explanations of experimental details that would disrupt the flow of the main text but nonetheless remain crucial to understanding and reproducing the research shown; figures of replicates for experiments of which representative data are shown in the main text can be added here if brief, or as Supplementary Data. Mathematical proofs of results not central to the paper can be added as an appendix.

\begin{table}[H] 
\caption{This is a table caption.\label{tab5}}
%\newcolumntype{C}{>{\centering\arraybackslash}X}
\begin{tabularx}{\textwidth}{CCC}
\toprule
\textbf{Title 1}	& \textbf{Title 2}	& \textbf{Title 3}\\
\midrule
Entry 1		& Data			& Data\\
Entry 2		& Data			& Data\\
\bottomrule
\end{tabularx}
\end{table}

\section[\appendixname~\thesection]{}
All appendix sections must be cited in the main text. In the appendices, Figures, Tables, etc. should be labeled, starting with ``A''---e.g., Figure A1, Figure A2, etc.

%%%%%%%%%%%%%%%%%%%%%%%%%%%%%%%%%%%%%%%%%%
\isPreprints{}{% This command is only used for ``preprints''.
\begin{adjustwidth}{-\extralength}{0cm}
} % If the paper is ``preprints'', please uncomment this parenthesis.
%\printendnotes[custom] % Un-comment to print a list of endnotes

\reftitle{References}

% Please provide either the correct journal abbreviation (e.g. according to the “List of Title Word Abbreviations” http://www.issn.org/services/online-services/access-to-the-ltwa/) or the full name of the journal.
% Citations and References in Supplementary files are permitted provided that they also appear in the reference list here. 

%=====================================
% References, variant A: external bibliography
%=====================================
% \bibliography{your_external_BibTeX_file}

%=====================================
% References, variant B: internal bibliography
%=====================================

% ACS format
\isAPAandChicago{}{%
\begin{thebibliography}{999}

\bibitem[Schick et~al.(2023)]{ref-toolformer}
Schick, T.; Dwivedi-Yu, J.; Dess\`{i}, R.; Raileanu, R.; Lomeli, M.; Zettlemoyer, L.; Cancedda, N.; Scialom, T.
Toolformer: Language Models Can Teach Themselves to Use Tools.
In \textit{Advances in Neural Information Processing Systems 36 (NeurIPS 2023)}; Curran Associates: Red Hook, NY, USA, 2023.

\bibitem[Yao et~al.(2023)]{ref-react}
Yao, S.; Zhao, J.; Yu, D.; Du, N.; Shafran, I.; Narasimhan, K.; Cao, Y.
ReAct: Synergizing Reasoning and Acting in Language Models.
In \textit{Proceedings of the 11th International Conference on Learning Representations (ICLR 2023)};
OpenReview.net, 2023.

\bibitem[Qin et~al.(2024)]{ref-toolllm}
Qin, Y.; Liang, S.; Ye, Y.; Zhu, K.; Yan, L.; Lu, Y.; Lin, Y.; Cong, X.; Tang, X.; Qian, B.; et~al.
ToolLLM: Facilitating Large Language Models to Master 16000+ Real-world APIs.
In \textit{Proceedings of the 12th International Conference on Learning Representations (ICLR 2024)};
OpenReview.net, 2024.

\bibitem[Patil et~al.(2023)]{ref-gorilla}
Patil, S.G.; Zhang, T.; Wang, X.; Gonzalez, J.E.
Gorilla: Large Language Model Connected with Massive APIs.
\textit{arXiv} \textbf{2023}, arXiv:2305.15334.

\bibitem[Ma et~al.(2024)]{ref-sciagent}
Ma, Z.; Gou, Z.; Hao, J.; Xu, K.; Hu, Z.; Yin, C.; Zheng, H.; Huang, M.; Deng, Y.
SciAgent: Tool-Augmented Language Models for Scientific Reasoning.
In \textit{Proceedings of the 2024 Conference on Empirical Methods in Natural Language Processing (EMNLP 2024)};
Association for Computational Linguistics: Stroudsburg, PA, USA, 2024.

\bibitem[Mialon et~al.(2023)]{ref-augmented-survey}
Mialon, G.; Dess\`{i}, R.; Lomeli, M.; Nalmpantis, C.; Pasunuru, R.; Raileanu, R.; Rozi\`{e}re, B.;
Schick, T.; Dwivedi-Yu, J.; Celikyilmaz, A.; et~al.
Augmented Language Models: A Survey.
\textit{Trans. Mach. Learn. Res.} \textbf{2023}.

\bibitem[Kim et~al.(2024)]{ref-llmcompiler}
Kim, S.; Moon, S.; Tabrizi, R.; Lee, N.; Mahoney, M.W.; Keutzer, K.; Gholami, A.
An LLM Compiler for Parallel Function Calling.
In \textit{Proceedings of the 41st International Conference on Machine Learning (ICML 2024)};
PMLR: Honolulu, HI, USA, 2024; Vol. 235, pp. 24167--24190.

\bibitem[Wang et~al.(2024)]{ref-codeact}
Wang, X.; Chen, Y.; Yuan, L.; Zhang, Y.; Li, Y.; Peng, H.; Ji, H.
Executable Code Actions Elicit Better LLM Agents.
In \textit{Proceedings of the 41st International Conference on Machine Learning (ICML 2024)};
PMLR: Honolulu, HI, USA, 2024; Vol. 235, pp. 50908--50933.

\bibitem[Nichols et~al.(2024)]{ref-pareval}
Nichols, D.; Davis, J.H.; Xie, Z.; Rajaram, A.; Bhatele, A.
Can Large Language Models Write Parallel Code?
In \textit{Proceedings of the 33rd International Symposium on High-Performance Parallel and Distributed Computing (HPDC 2024)};
ACM: New York, NY, USA, 2024; pp. 117--130.

\bibitem[Nichols et~al.(2024b)]{ref-hpccoder}
Nichols, D.; Davis, J.H.; Bhatele, A.
HPC-Coder-V2: Studying Code LLMs Across Low-Resource Parallel Languages.
\textit{arXiv} \textbf{2024}, arXiv:2412.15178.

\bibitem[Xu et~al.(2025)]{ref-deepseekhpc}
Xu, C.; Li, M.; Xu, X.; Qin, X.; Chen, P.
LLM \& HPC: Benchmarking DeepSeek's Performance in High-Performance Computing Tasks.
In \textit{Proceedings of ISC High Performance 2025};
Springer: Cham, Switzerland, 2025.

\bibitem[He et~al.(2025)]{ref-legionllm}
He, W.; Shen, Z.; Nie, A.; Luo, C.; Nus, T.; Aiken, A.; Zaharia, M.
Improving Parallel Program Performance with LLM Optimizers via Agent-System Interfaces.
In \textit{Proceedings of the 42nd International Conference on Machine Learning (ICML 2025)};
PMLR, 2025.

\bibitem[Mialon et~al.(2024)]{ref-gaia}
Mialon, G.; Fourrier, C.; Swift, C.; Wolf, T.; LeCun, Y.; Scialom, T.
GAIA: A Benchmark for General AI Assistants.
In \textit{Proceedings of the 12th International Conference on Learning Representations (ICLR 2024)};
OpenReview.net, 2024.

\bibitem[Jim\'{e}nez et~al.(2024)]{ref-swebench}
Jim\'{e}nez, C.E.; Yang, J.; Wettig, A.; Yao, S.; Pei, K.; Press, O.; Narasimhan, K.
SWE-bench: Can Language Models Resolve Real-World GitHub Issues?
In \textit{Proceedings of the 12th International Conference on Learning Representations (ICLR 2024)};
OpenReview.net, 2024.

\bibitem[Chen et~al.(2025)]{ref-sciagentbench}
Chen, Z.; Shi, Z.; Wang, S.; et~al.
ScienceAgentBench: Toward Rigorous Assessment of Language Agents for Data-Driven Scientific Discovery.
In \textit{Proceedings of the 13th International Conference on Learning Representations (ICLR 2025)};
OpenReview.net, 2025.

\bibitem[Xu et~al.(2024)]{ref-agentcompany}
Xu, F.F.; Ye, Y.; Jain, A.; Du, Y.; et~al.
TheAgentCompany: Benchmarking LLM Agents on Consequential Real World Tasks.
\textit{arXiv} \textbf{2024}, arXiv:2412.14161.

\bibitem[Anthropic(2024)]{ref-mcp}
Anthropic.
Model Context Protocol.
Available online: \url{https://modelcontextprotocol.io} (accessed on 1 May 2026).

\bibitem[Hou et~al.(2025)]{ref-mcpsurvey}
Hou, X.; Zhao, Y.; Wang, S.; Wang, H.
Model Context Protocol (MCP): Landscape, Security Threats and Future Research Directions.
\textit{arXiv} \textbf{2025}, arXiv:2503.23278.

\bibitem[Luo et~al.(2025)]{ref-mcpsmell}
Luo, L.; et~al.
Model Context Protocol (MCP) Tool Descriptions Are Smelly!
\textit{arXiv} \textbf{2025}, arXiv:2602.14878.

\end{thebibliography}
}

% Chicago format (Used for journal: arts, genealogy, histories, humanities, jintelligence, laws, literature, religions, risks, socsci)
\isChicagoStyle{%
\begin{thebibliography}{999}
% Reference 1
\bibitem[Aranceta-Bartrina(1999a)]{ref-journal}
Aranceta-Bartrina, Javier. 1999a. Title of the cited article. \textit{Journal Title} 6: 100--10.
% Reference 2
\bibitem[Aranceta-Bartrina(1999b)]{ref-book1}
Aranceta-Bartrina, Javier. 1999b. Title of the chapter. In \textit{Book Title}, 2nd ed. Edited by Editor 1 and Editor 2. Publication place: Publisher, vol. 3, pp. 54–96.
% Reference 3
\bibitem[Baranwal and Munteanu {[1921]}(1955)]{ref-book2}
Baranwal, Ajay K., and Costea Munteanu. 1955. \textit{Book Title}. Publication place: Publisher, pp. 154--96. First published 1921 (op-tional).
% Reference 4
\bibitem[Berry and Smith(1999)]{ref-thesis}
Berry, Evan, and Amy M. Smith. 1999. Title of Thesis. Level of Thesis, Degree-Granting University, City, Country. Identifi-cation information (if available).
% Reference 5
\bibitem[Cojocaru et al.(1999)]{ref-unpublish}
Cojocaru, Ludmila, Dragos Constatin Sanda, and Eun Kyeong Yun. 1999. Title of Unpublished Work. \textit{Journal Title}, phrase indicating stage of publication.
% Reference 6
\bibitem[Driver et al.(2000)]{ref-proceeding}
Driver, John P., Steffen Rohrs, and Sean Meighoo. 2000. Title of Presentation. In \textit{Title of the Collected Work} (if available). Paper presented at Name of the Conference, Location of Conference, Date of Conference.
% Reference 7
\bibitem[Harwood(2008)]{ref-url}
Harwood, John. 2008. Title of the cited article. Available online: URL (accessed on Day Month Year).
\end{thebibliography}
}{}

% APA format (Used for journal: admsci, behavsci, businesses, econometrics, economies, education, ejihpe, games, humans, ijfs, journalmedia, jrfm, languages, psycholint, publications, tourismhosp, youth)
\isAPAStyle{%
\begin{thebibliography}{999}
% Reference 1
\bibitem[\protect\citeauthoryear{Azikiwe \BBA\ Bello}{{2020a}}]{ref-journal}
Azikiwe, H., \& Bello, A. (2020a). Title of the cited article. \textit{Journal Title}, \textit{Volume}(Issue), 
Firstpage--Lastpage/Article Number.
% Reference 2
\bibitem[\protect\citeauthoryear{Azikiwe \BBA\ Bello}{{2020b}}]{ref-book1}
Azikiwe, H., \& Bello, A. (2020b). \textit{Book title}. Publisher Name.
% Reference 3
\bibitem[Davison(1623/2019)]{ref-book2}
Davison, T. E. (2019). Title of the book chapter. In A. A. Editor (Ed.), \textit{Title of the book: Subtitle} 
(pp. Firstpage--Lastpage). Publisher Name. (Original work published 1623) (Optional).
% Reference 4
\bibitem[Fistek et al.(2017)]{ref-proceeding}
Fistek, A., Jester, E., \& Sonnenberg, K. (2017, Month Day). Title of contribution [Type of contribution]. Conference Name, Conference City, Conference Country.
% Reference 5
\bibitem[Hutcheson(2012)]{ref-thesis}
Hutcheson, V. H. (2012). \textit{Title of the thesis} [XX Thesis, Name of Institution Awarding the Degree].
% Reference 6
\bibitem[Lippincott \& Poindexter(2019)]{ref-unpublish}
Lippincott, T., \& Poindexter, E. K. (2019). \textit{Title of the unpublished manuscript} [Unpublished manuscript/Manuscript in prepara-tion/Manuscript submitted for publication]. Department Name, Institution Name.
% Reference 7
\bibitem[Harwood(2008)]{ref-url}
Harwood, J. (2008). \textit{Title of the cited article}. Available online: URL (accessed on Day Month Year).
\end{thebibliography}
}{}

% If authors have biography, please use the format below
%\section*{Short Biography of Authors}
%\bio
%{\raisebox{-0.35cm}{\includegraphics[width=3.5cm,height=5.3cm,clip,keepaspectratio]{Definitions/author1.pdf}}}
%{\textbf{Firstname Lastname} Biography of first author}
%
%\bio
%{\raisebox{-0.35cm}{\includegraphics[width=3.5cm,height=5.3cm,clip,keepaspectratio]{Definitions/author2.jpg}}}
%{\textbf{Firstname Lastname} Biography of second author}

% For the MDPI journals use author-date citation, please follow the formatting guidelines on http://www.mdpi.com/authors/references
% To cite two works by the same author: \citeauthor{ref-journal-1a} (\citeyear{ref-journal-1a}, \citeyear{ref-journal-1b}). This produces: Whittaker (1967, 1975)
% To cite two works by the same author with specific pages: \citeauthor{ref-journal-3a} (\citeyear{ref-journal-3a}, p. 328; \citeyear{ref-journal-3b}, p.475). This produces: Wong (1999, p. 328; 2000, p. 475)

%%%%%%%%%%%%%%%%%%%%%%%%%%%%%%%%%%%%%%%%%%
%% for journal Sci
%\reviewreports{\\
%Reviewer 1 comments and authors’ response\\
%Reviewer 2 comments and authors’ response\\
%Reviewer 3 comments and authors’ response
%}
%%%%%%%%%%%%%%%%%%%%%%%%%%%%%%%%%%%%%%%%%%
\PublishersNote{}
\isPreprints{}{% This command is only used for ``preprints''.
\end{adjustwidth}
} % If the paper is ``preprints'', please uncomment this parenthesis.
\end{document}

